% TOP_PROBLEM: {"id": 30001019, "title": "Algebraic Dimension and Reduction of Hyperkähler Manifolds", "category_id": 5, "category": {"id": 5, "name": "algebraic_geometry", "display_name": "Algebraic Geometry", "description": "Geometric objects defined by polynomial equations.", "slug": "algebraic-geometry", "order_index": 5, "created_at": "2026-07-31T15:26:25.671Z"}, "status": "partially_solved"}
% FIRST_POSTED: 2026-10-06
% ATTEMPT_STATUS: unresolved
% OPENAI_CLAIM: {"schema_version":1,"source_repo":"https://github.com/openai/math","source_commit":"adc7f1241b42e322a6451854ab7e4b4c146bf78a","release_date":"2026-10-06","match":"partial","resolution":"partial","independently_reviewed":false,"second_pass":"n/a","families":[{"family":"041","title":"Hyperkähler SYZ and projective-space bases","subject":"Algebraic and complex geometry","manuscripts":[{"title":"The strong hyperkähler SYZ conjecture","pdf_path":"preprints/The-Strong-Hyperkahler-SYZ-Conjecture-September-23-2026/main.pdf","readme_path":"preprints/The-Strong-Hyperkahler-SYZ-Conjecture-September-23-2026/README.md","date":"2026-09-23","theorem_ref":"thm:main (sections/introduction.tex:25); cor:finiteness (sections/introduction.tex:57)","statement_sources":[{"path":"preprints/The-Strong-Hyperkahler-SYZ-Conjecture-September-23-2026/build/sections/introduction.tex","line":25,"label":"thm:main"},{"path":"preprints/The-Strong-Hyperkahler-SYZ-Conjecture-September-23-2026/build/sections/introduction.tex","line":57,"label":"cor:finiteness"}]}],"lean":null,"reasoning_trace":null}]}
\documentclass[11pt]{article}
\usepackage[margin=25mm]{geometry}
\usepackage{fontspec}
\setmainfont{Latin Modern Roman}
\usepackage{amsmath,amssymb,mathtools}
\usepackage{unicode-math}
\setmathfont{Latin Modern Math}
\usepackage{enumitem}
\usepackage{xurl,hyperref,bookmark}
\hypersetup{colorlinks=true,urlcolor=blue}
\setcounter{secnumdepth}{0}
\setlength{\parindent}{0pt}
\setlength{\parskip}{5pt}
\setlength{\emergencystretch}{3em}
\begin{document}
\section*{Algebraic Dimension and Reduction of Hyperkähler Manifolds}

\subsection{OpenAI's claim}
On a compact Kähler IHS manifold, a holomorphic line is claimed to be semiample whenever its Chern class is nonzero, nef, and isotropic, regardless of deformation type. The cited manuscript argues for that result; its relation to this problem is specified below.

\subsection{Mathematical statement of the claimed result}
\paragraph{Mathematical claim 1.} Let $X$ be an irreducible holomorphic symplectic manifold, required to be compact and Kähler, of complex dimension $2n$, $n\ge1$. Write $q_X$ for its Beauville--Bogomolov form and $\overline{\mathcal K_X}$ for the closed Kähler cone. For a holomorphic line bundle $L$ with nonzero $c_1(L)\in\overline{\mathcal K_X}$ and $q_X(c_1(L))=0$, OpenAI claims semiampleness and, more precisely,
\[
L^m\simeq f^*A,\qquad f:X\twoheadrightarrow B,\qquad\dim_{\mathbb C}B=n,
\]
with $m>0$, connected fibers, normal projective $B$ and ample $A$. Every fiber component has dimension $n$ and the holomorphic symplectic form vanishes on its smooth locus. A separate corollary claims finitely many smooth deformation types in each dimension $2n$ among these manifolds with $b_2\ge5$.

\subsection{Outline of the argument}
Manuscript 1 (family 041, 2026-09-23): The prescribed lattice vector is retained through a degeneration that produces a special Lagrangian torus. Rotation makes that torus holomorphic on a nonprojective deformation, where the torus-fibration theorem supplies semiampleness; the fixed-class deformation result brings the conclusion back to the original manifold.

\subsection{Scope relative to this problem}
Manuscript 1 (family 041, 2026-09-23): A nonzero nef isotropic line produces a Lagrangian fibration after semiampleness, providing the expected intermediate algebraic reduction in that additional-line class. The record asks the algebraic-dimension classification for arbitrary hyperkahler manifolds without this nef-line hypothesis.

\subsection{Formal verification}
Family 041 has no Lean scope page in this snapshot; no formalization of these manuscript results is listed.

This site has not executed the Lean files or checked their dependencies. These descriptions report the scope of the linked declarations, rather than a new verification of the proofs.

\subsection{Status caveat}
OpenAI attributes these results to an internal model. Its release notice cautions that results lacking full formalization can contain mistakes. This site has not independently assessed the proofs; a solved label records the scope of the reported claim.

\subsection{Sources}
\begin{itemize}
\item OpenAI, \emph{The strong hyperkähler SYZ conjecture}, 2026-09-23; citation key \texttt{OAI:The-Strong-Hyperkahler-SYZ-Conjecture-September-23-2026}. \href{https://github.com/openai/math/blob/adc7f1241b42e322a6451854ab7e4b4c146bf78a/preprints/The-Strong-Hyperkahler-SYZ-Conjecture-September-23-2026/main.pdf}{PDF (GitHub)}; \href{https://github.com/openai/math/raw/adc7f1241b42e322a6451854ab7e4b4c146bf78a/preprints/The-Strong-Hyperkahler-SYZ-Conjecture-September-23-2026/main.pdf}{PDF (direct)}; \href{https://github.com/openai/math/blob/adc7f1241b42e322a6451854ab7e4b4c146bf78a/preprints/The-Strong-Hyperkahler-SYZ-Conjecture-September-23-2026/README.md}{Manuscript page and citation}; \href{https://github.com/openai/math/blob/main/preprints/The-Strong-Hyperkahler-SYZ-Conjecture-September-23-2026/main.pdf}{Latest version}; Statement source: \href{https://github.com/openai/math/blob/adc7f1241b42e322a6451854ab7e4b4c146bf78a/preprints/The-Strong-Hyperkahler-SYZ-Conjecture-September-23-2026/build/sections/introduction.tex\#L25}{thm:main}; Statement source: \href{https://github.com/openai/math/blob/adc7f1241b42e322a6451854ab7e4b4c146bf78a/preprints/The-Strong-Hyperkahler-SYZ-Conjecture-September-23-2026/build/sections/introduction.tex\#L57}{cor:finiteness}.
\item \href{https://github.com/openai/math/blob/adc7f1241b42e322a6451854ab7e4b4c146bf78a/CONTENTS.md}{OpenAI catalogue entry}.
\end{itemize}
\end{document}
